\documentclass[11pt,letterpaper]{article}
\usepackage{../styles/mae2250-handout}

\renewcommand{\HandoutNumber}{4}
\renewcommand{\HandoutTitle}{Coffee Grinder}
\renewcommand{\PartTitle}{Part 4: Final Steps}
\renewcommand{\DueDate}{}
\renewcommand{\TeamType}{Individual}

\begin{document}
\MakeHandoutHeader

Based on what you have learned in previous parts, you may alter your design and finalize it
by putting it in the housing. You must obey all assignment constraints.

\PartHeading{PART 4}
\begin{AssignmentSteps}
  \item \underline{CAD the final prototype.}
    \begin{itemize}
      \item Make a new assembly and move your internal mechanism into it.
        \begin{itemize}
          \item Make any changes as necessary
        \end{itemize}
      \item Model the final housing.
        \begin{itemize}
          \item Use your center of gravity test to ensure nothing tips over.
        \end{itemize}
    \end{itemize}
  \item \underline{Check for interfering parts.}
    \begin{itemize}
      \item Navigate to Inspect$\rightarrow$Section Analysis. Select a plain that meaningfully intersects your
        model, then drag it through your model to see if any of the internal parts intersect. Take a
        screenshot of your analysis where you can see the greatest cross section of your model.
      \item If you think you see an interfering part, navigate to Inspect$\rightarrow$Interference. Click on the two
        parts that you think interfere, then hit ``Compute.''
    \end{itemize}
  \item \underline{Writeup}
    \begin{itemize}
      \item In two brief paragraphs, describe the choices you made in this assignment as well as lessons
        learned. As always, we are interested in your thought process, not your writing prowess.
        \begin{itemize}
          \item In paragraph 1, focus on your design choices as well as anything that was difficult or
            interesting.
          \item In paragraph 2, focus on the feedback you received and the lessons you learned once
            you built your model. Note any revisions you made afterwards.
        \end{itemize}
    \end{itemize}
\end{AssignmentSteps}
\textbf{Ensure that your model is in the folder ``coffeegrinder\_Final,'' and submit a document with
your screenshot, writeup, and updated bill of materials to Canvas by \{part 4 due date\}.}

Now you have completed all of your individual assignments for the class. We hope that you use the
skills that you have learned in constructing your prototype for the ODP.

\newpage
\PartHeading{TA CHECKLIST}
\begin{Checklist}
  \item All pieces in the bill of materials match the CAD.
  \item The bill of materials does not exceed \$100.
  \item A screenshot of the section analysis was uploaded.
  \item A pair of gears and a conical burr was used in the design.
  \item The gears are properly motion-linked.
  \item All parts are properly constrained.
  \item The dimensions of the CAD do not exceed 300mm*200mm*300mm.
  \item The CAD follows good practices.
  \item Consideration of feedback is evident, either in the CAD model or in the writeup. *Note that
    not every piece of feedback given must be considered or implemented to receive the point.
  \item The writeup demonstrates the student's thought process and lessons learned.
\end{Checklist}

\textbf{\{At bare minimum, TAs should give feedback for parts 1 and 4.\}}

\end{document}
